\section{Subgroup and severity analysis}
\label{app:subgroup}

\begin{figure}[h]
\centering
\includegraphics[width=\linewidth]{fig_fairness.png}
\caption{\textbf{Left:} race-stratified AUC at epoch 100 with bootstrap
intervals. The black subgroup is lowest under every policy, but note that
\emph{every} group improves from \textsc{random} to \textsc{intensity}.
\textbf{Right:} max--min race gap against aggregate AUC across \NprobesSub{}
checkpoints. The fitted slope is positive but does not survive multiplicity
correction or branch-level aggregation (Table~\ref{tab:subtrends}), so we draw
no trend conclusion from it.}
\label{fig:fair}
\end{figure}

All numbers below come from saved test predictions joined to FairVision's
released metadata in deterministic loader order; the join is validated by exact
label reconstruction (\Ntest/\Ntest{} labels recovered in every probe). The
\NprobesSub{} probes used here are exactly those the inventory marks valid: the
four retracted coverage probes of Appendix~\ref{app:excluded} and the two
precision-spliced \textsc{anatomy-v2} probes are excluded here as everywhere
else in the paper.

\paragraph{Why we report two sample sizes.} The \NprobesSub{} probes span only
\NbranchesSub{} pretraining branches, because several are different epochs of
one training run. Probes within a branch share an encoder lineage, a seed and
the test split, so treating them as \NprobesSub{} independent observations is
pseudo-replication. Table~\ref{tab:subtrends} therefore reports each trend twice:
once per checkpoint, and once after averaging within each branch. Where the two
disagree, the branch-level figure is the one to believe.

\begin{table}[h]
\centering
\caption{Subgroup gap trends. ``gap range'' is the max--min subgroup AUC gap
observed across the \NprobesSub{} valid probes. $q$ is Benjamini--Hochberg
across the seven attributes tested, which is necessary because seven trends were
examined. Only the sex gap and the severity gap survive correction, and both
\emph{narrow} as aggregate AUC rises.}
\label{tab:subtrends}
\footnotesize
\setlength{\tabcolsep}{4pt}
\begin{tabular}{llccccccc}
\toprule
& & & \multicolumn{3}{c}{per checkpoint ($n{=}27$)} & \multicolumn{2}{c}{per branch ($n{=}11$)} \\
\cmidrule(lr){4-6}\cmidrule(lr){7-8}
attribute & worst group & gap range & $\rho$ & $p$ & $q$ & $\rho$ & $p$ \\
\midrule
sex & female (27/27) & 0.0270--0.0396 & -0.739 & $<$0.0001 & $<$0.0001 & -0.891 & 0.0002 \\
race & black (27/27) & 0.0433--0.0935 & +0.487 & 0.0100 & 0.0232 & +0.518 & 0.1025 \\
ethnicity & hispanic (27/27) & 0.0203--0.0661 & +0.151 & 0.4509 & 0.4509 & +0.409 & 0.2115 \\
language & other languages (22/27) & 0.0518--0.0805 & +0.184 & 0.3589 & 0.4187 & +0.400 & 0.2229 \\
marital status & single (27/27) & 0.0349--0.0654 & -0.380 & 0.0503 & 0.0881 & -0.373 & 0.2589 \\
age & $<$60 (26/27) & 0.0474--0.0704 & +0.300 & 0.1279 & 0.1791 & +0.236 & 0.4841 \\
disease severity & mild ($-6$,$-2$] (27/27) & 0.1266--0.1400 & -0.551 & 0.0029 & 0.0102 & -0.618 & 0.0426 \\
\bottomrule
\end{tabular}

\end{table}

\paragraph{Race.} The black subgroup has the lowest AUC in every one of the
\NprobesSub{} probes ($n{=}\NBlack$ black, $n{=}\NAsian$ asian, $n{=}\NWhite$
white). At matched epoch 100 the black subgroup improves from
$\RaceRandomBlack$ under \textsc{random} to $\RaceOracleBlack$ under
\textsc{intensity} ($\BlackGainOracle$), a larger absolute gain than the white
subgroup's ($\RaceRandomWhite \to \RaceOracleWhite$). The max--min gap
nonetheless widens slightly, because the asian subgroup --- already the
highest --- gains most ($\RaceRandomAsian \to \RaceOracleAsian$). A widening
max--min gap with universally rising subgroup AUCs is not evidence of harm, and
we do not present it as such. The checkpoint-level correlation between aggregate
AUC and racial gap is $\rho{=}\SubRaceRho$ ($p{=}\SubRaceRhoP$), which fails BH
correction ($q{=}\SubRaceQ$) and vanishes at branch level
($\rho{=}\SubRaceBranchRho$, $p{=}\SubRaceBranchP$). We therefore explicitly do
\emph{not} claim that more accurate models are less fair.

\paragraph{Disease severity.} FairVision's label is defined by visual-field mean
deviation (md): every test volume with md $\leq -2$ is positive. Severity
therefore cannot be used as an ordinary subgroup axis, because within-bin AUC is
undefined --- a trap that silently yields nothing under naive stratification.
Instead each positive stratum is scored against the shared pool of all
\Nneg{} negatives.

\begin{table}[h]
\centering
\caption{Severity-stratified AUC at matched epoch 100. Each positive stratum is
scored against the shared pool of all \Nneg{} negatives. Every stratum improves,
and mild disease improves most.}
\label{tab:severity}
\small
\begin{tabular}{lccccc}
\toprule
stratum & $n_{+}$ & \textsc{random} & \textsc{envelope} & \textsc{intensity} & $\Delta$ \\
\midrule
severe (md $\leq -12$)          & 334 & \SevRandomSevere   & \SevEnvelopeSevere   & \SevOracleSevere   & \SevDeltaSevere \\
moderate ($-12 <$ md $\leq -6$) & 460 & \SevRandomModerate & \SevEnvelopeModerate & \SevOracleModerate & \SevDeltaModerate \\
mild ($-6 <$ md $\leq -2$)      & 672 & \SevRandomMild     & \SevEnvelopeMild     & \SevOracleMild     & \textbf{\SevDeltaMild} \\
\bottomrule
\end{tabular}
\end{table}

The severe-to-mild gap is $\SevGapRandom$ for the null and $\SevGapOracle$ for
the strongest policy, and stays within $\SeverityGapSpread$
($\SubSeverityGapMin$--$\SubSeverityGapMax$) across all \NprobesSub{} probes
while aggregate AUC spans $\SubAUCMin$--$\SubAUCMax$. Target placement raises
the whole severity curve --- mild disease most of all --- without altering the
ordering that makes early disease hard. We note that mild functional loss is
harder to separate from normal OCT by construction, so this ordering is expected
clinically; what the data add is that no masking policy we tested changes it.

\paragraph{Limitations.} One dataset, one test split, reused throughout this
programme's history. Because policies, checkpoints and analyses were chosen
after repeated inspection of that split, the $p$-values here should be read as
descriptive rather than as confirmatory inference. Race and ethnicity are
self-reported categories inherited from FairVision's coding; at these subgroup
sizes we cannot speak to intersectional groups. All figures are frozen mean-pool
linear probes with AUC as the only metric; we report no subgroup calibration,
no sensitivity at a fixed specificity, and no threshold-transfer analysis, so
this audit should not be read as a complete fairness evaluation. Disparities
could differ under fine-tuning or a stronger head.

